\chapter{Requirements, specifications, and stakeholders}

\section{From requests to obligations}
“Make checkout faster” is a request, not yet a requirement. It does not say
which users, which journey, which measurement boundary, or what trade-off is
acceptable. A useful requirement names an actor, an observable behaviour, a
condition, and a threshold or invariant when one matters.

One Ledgerline requirement can be written as:

\begin{quote}
When a caller submits a syntactically valid order with an unused idempotency key,
the service shall persist one accepted order and return its identifier. Repeating
the same request with the same key shall return the original result without
creating another order.
\end{quote}

This statement creates questions: How long is a key retained? What if the first
request timed out after commit? Must the payload match on reuse? What error does a
conflicting payload receive? Requirements become useful when their boundary
cases are discussed rather than hidden.

\section{Stakeholder analysis}
A stakeholder is anyone who can affect, use, constrain, be affected by, or be
held accountable for the system. The person requesting a feature is not always
the person who bears its risks. A small table is often enough to uncover this.

\noindent\begin{tabularx}{\textwidth}{@{}p{0.19\textwidth}p{0.24\textwidth}X@{}}
\toprule
Stakeholder & Interest & Questions \\
\midrule
Customer & A correct, understandable order & What counts as confirmation? Can the order be retried safely? \\
Support & Ability to explain outcomes & Can a person find the request, status, and reason for rejection? \\
Operator & Stable service and recoverable data & What is monitored, backed up, and safe to change? \\
Finance & Reconciled money & Which values are authoritative, and how are corrections audited? \\
Security and privacy & Limited exposure and accountability & Which data is collected, for what purpose, and under whose authority? \\
\bottomrule
\end{tabularx}

Stakeholder analysis does not mean accepting every preference. It makes the
conflict explicit. A customer may want unlimited retention for convenience while
privacy policy requires deletion. An operator may want a safe pause while a
commercial team wants availability. Engineering supplies options and evidence;
governance decides which values prevail.

\section{Functional and quality requirements}
Functional requirements describe behaviour: create an order, cancel it, export a
record. Quality requirements describe conditions on that behaviour: a latency
percentile, a recovery time, an accessibility criterion, or a permitted data
residency. The categories interact. “Returns a result” is incomplete if the
result is too late for the user, cannot be accessed by keyboard, or exposes a
secret in logs.

Use a requirement record with at least:

\begin{enumerate}
  \item identifier and short name;
  \item rationale and stakeholders;
  \item normative statement;
  \item acceptance examples, including invalid cases;
  \item measurable quality target and measurement boundary;
  \item dependencies, assumptions, and unresolved decisions;
  \item owner and review date.
\end{enumerate}

\section{Specifications as executable agreements}
A specification is valuable when independent people can use it to predict the
same behaviour. Examples include an OpenAPI document, a SQL schema, a property
test, a protocol state machine, or a runbook. A prose requirement remains
important for purpose and context, but a machine-checkable representation catches
drift earlier.

Do not make the specification larger than the uncertainty it resolves. A stable
public API deserves explicit compatibility rules. A private function may need
only a type, a test, and a clear name. Excess ceremony creates documents that
are neither read nor updated.

\section{Acceptance and examples}
Examples expose ambiguity. The following set describes the same rule from three
angles:

\begin{itemize}
  \item Given two units at 1,250 cents, the subtotal is 2,500 cents.
  \item Given quantity zero, the request is rejected before persistence.
  \item Given the same idempotency key twice with the same payload, the second
        response refers to the first order.
\end{itemize}

Examples are not a replacement for general rules. They are probes that make the
rule concrete and reveal missing decisions. Property-based tests can then state
the broader invariant: for every valid collection of positive integer prices and
quantities, the total is an integer and equals the sum of each product.

\section{Change and traceability}
Requirements change because the world changes. A useful traceability chain is:
requirement -> design decision -> code -> automated check -> operational signal.
The chain should be lightweight enough to maintain. A commit message or issue
identifier can be sufficient for a small team; regulated work may need a formal
record. The principle is the same: a future reader should be able to ask why a
constraint exists and find evidence.

\begin{tradeoffs}
Detailed upfront specifications reduce some ambiguity but can freeze wrong
assumptions. Pure improvisation learns quickly but can leave invisible contracts.
Use more formality where the cost of misunderstanding is high, and keep the
feedback cycle short where uncertainty is high.
\end{tradeoffs}

\begin{exercise}
Turn “the report page should feel responsive” into a measurable requirement.
State the user journey, the measurement boundary, the percentile and time window,
the test data shape, and one reason the target might need revision.
\end{exercise}

\section{Chapter summary}
Requirements become engineering material when they identify actors, behaviour,
conditions, and evidence. Include quality properties, examine stakeholders who
bear risks, use examples to expose ambiguity, and keep a trace from intent to
operation without creating paperwork that nobody can maintain.
